\section{Discrete equivalent layers on a plane}
\label{eq:planar_eqlayer}

In practical geophysical applications involving discrete field data, 
the interpreter never use a continuous equivalent layer $p_{\varepsilon}$ 
(Eq. \ref{eq:p-Fourier}).
Instead, one typically estimates a ``discrete equivalent layer'' that is 
sufficiently smooth and fits the observed  discrete data within a prescribed error 
tolerance \cite{oliveirajr_etal2026_1}. 
We regard any discrete layer commonly obtained in practice  
as an approximation of the corresponding continuous equivalent layer $p_{\varepsilon}$ 
defined by using the same differential operators $\alpha_{d}$ and $\alpha_{a}$
(Eq(s). \ref{eq:diff-operators-generic} and \ref{eq:p-Fourier}). 
A ``discrete equivalent layer'' on the horizontal surface  
$\mathcal{S}''$ (Fig.  \ref{fig:Figure1}) can be formally defined as follows:  

\begin{definition}[Discrete equivalent layer]
	Let us assume to have an $N \times 1$ observed data vector $\mathbf{d}^{o}$ whose elements 
	are the potential-field data $d^{o}_{i}$ at observation points $\mathbf{r}_{i}$ outside the sources, 
	$i \in \{1, \dots, N\}$, above a horizontal surface $\mathcal{S}''$ with constant vertical 
	coordinate $z''_{0}$. Besides, consider an $N \times M$ 
	matrix $\mathbf{A}$ with element $a_{ij}$ defined by a harmonic function that depends only on $\mathbf{r}_{i}$ 
	and $\mathbf{r}''_{j} \in \mathcal{S}''$, $j \in \{1, \dots, M\}$, and
	vanishes as the distance $\| \mathbf{r}_{i} - \mathbf{r}''_{j} \|$ increases. 
	Then, for a chosen accuracy $\epsilon > 0$, we may call the $M \times 1$ vector $\mathbf{p}_{\epsilon}$ with $j$-th element $p_{j}$ 
	a \underline{discrete equivalent layer} if its predicted potential field 
	\begin{equation}
		\mathbf{d}_{\epsilon} = \mathbf{A} \, \mathbf{p}_{\epsilon}
		\label{eq:predicted-data-discrete-layer}
	\end{equation}
	satisfies 
	\begin{equation}
		\| \mathbf{d}^{o} - \mathbf{d}_{\epsilon} \|_{2}^{2} \le \epsilon \: .
		\label{eq:epsilon-bound-data-discrete}
	\end{equation}
	\label{def:discrete-eqlayer}
\end{definition}

In Definition \ref{def:discrete-eqlayer}, $\mathbf{d}_{\epsilon}$ is the predicted data vector 
whose $i$-th element represents the potential field due to $\mathbf{p}_{\epsilon}$ at the 
observation point $\mathbf{r}_{i}$. 
Besides, the elements of $\mathbf{A}$ are usually defined by 
$a_{ij} \equiv D_{a} \Psi(\mathbf{r}_{i}, \mathbf{r}''_{j})$, which represents a 
partial derivative of $\Psi$ (Eq. \ref{eq:Psi}), with $D_{a}$ given by the space-domain 
differential operator $D$ (Eq. \ref{eq:diff-operators-generic}). 
It is important to bear in mind that the horizontal surface $\mathcal{S}''$ in Definition 
\ref{def:discrete-eqlayer} may cross the sources.
The main assumption in Definition \ref{def:discrete-eqlayer} is that $d^{o}_{i}$ can be 
approximated within a given accuracy $\epsilon$ by the potential field produced by a discrete 
and finite set of $M$ virtual sources such as monopoles (or point masses) or dipoles with 
constant and unknown physical property $p_{j}$ (e.g., mass, magnetic-moment intensity). 
These virtual sources are usually called ``equivalent sources''.
We emphasize that Definition \ref{def:discrete-eqlayer} presumes a planar equivalent layer 
formed by punctual equivalent sources that satisfies the data fit criterion given by Eq. 
\ref{eq:epsilon-bound-data-discrete}, but these premises are certainly open to adjustments if 
necessary.

After estimating a particular $\mathbf{p}_{\epsilon}$ (Definition \ref{def:discrete-eqlayer}), 
a potential field transformation such as upward continuation, 
interpolation or reduction to the pole, for example, can be computed as follows:
\begin{equation}
	\mathbf{f}_{\epsilon} = \mathbf{T} \, \mathbf{p}_{\epsilon} \: ,
	\label{eq:transformed-field-discrete-layer} 
\end{equation}
where $\mathbf{f}_{\epsilon}$ is an $L \times 1$ vector whose $k$-th element approximates 
the transformed potential field $f(\mathbf{r}^{\ddagger}_{k})$ 
(Eq. \ref{eq:transformed-field-space}) 
at a point $\mathbf{r}^{\ddagger}_{k} \in \mathcal{S}^{\ddagger}$ above $\mathcal{S}''$ 
and $\mathbf{T}$ is an $L \times M$ matrix whose element 
$t_{kj} \equiv D_{t} \Psi(\mathbf{r}^{\ddagger}_{k}, \mathbf{r}''_{j})$,
$k \in \{ 1, \dots, L \}$, $j \in \{ 1, \dots, M \}$, 
is commonly defined by a partial derivative of $\Psi$ (Eq. \ref{eq:Psi}), with $D_{t}$ 
given by the space-domain differential operator $D$ (Eq. \ref{eq:diff-operators-generic}).
The element $t_{kj}$ represents the transformed potential field due to the $j$-th equivalent 
source. 
We have defined the transformed potential field at points 
$\mathbf{r}^{\ddagger}_{k} \in \mathcal{S}^{\ddagger}$ above $\mathcal{S}''$ , however they 
can be defined at any place above $\mathcal{S}''$, including the observation points. 

From the physical point of view, the product of the $j$-th element $p_{j}$ of 
$\mathbf{p}_{\epsilon}$ and the harmonic function $a_{ij}$ representing the element $ij$ of 
matrix $\mathbf{A}$ (Eq. \ref{eq:predicted-data-discrete-layer}) defines the potential 
field produced at the $i$-th observation point by the $j$-th equivalent source with scalar 
physical property $p_{j}$.
For example, if $\mathbf{d}^{o}$ contains gravity disturbance data, the differential operator 
$D_{a}$ is commonly set as $D(1, \hat{\mathbf{z}}) = \tfrac{\partial}{\partial z}$ 
(Eq. \ref{eq:diff-operators-generic}), where $\hat{\mathbf{z}} = (0, 0, 1)^{\top}$, so that 
the term $p_{j} \, a_{ij}$ defines the vertical component of the gravitational attraction
produced at $\mathbf{r}_{i}$ by a monopole located at
$\mathbf{r}''_{j} \in \mathcal{S}''$, with mass $p_{j}$. 
This is the base of the pioneer approach presented by \citep{dampney1969} for processing gravity data.
For the case in which $\mathbf{d}^{o}$ contains total-field anomaly data 
\citep[e.g.,][p. 178]{blakely1996} instead, the common approach consists in
following \citep{emilia1973} and setting $D_{a} \equiv D(1, \hat{\mathbf{u}}) \, D(1, \hat{\mathbf{v}})$, 
where $\hat{\mathbf{u}}$ and $\hat{\mathbf{h}}$ define the directions of, respectively, the local main field and 
the total magnetization of the true sources, so that the term $p_{j} \, a_{ij}$ represents the approximated 
total-field anomaly \citep[e.g.,][p. 179]{blakely1996} produced at the $i$-th observation point by a dipole 
located at $\mathbf{r}''_{j} \in \mathcal{S}''$, with magnetic moment intensity $p_{j}$.